\section{Classification complète des particules et des états physiques}
\label{sec:catalogo-realizaciones-tipadas-rev8}
\index{secteurs physiques!inventaire typé}
\index{atlas!81/56/13/324}

Les réalisations sont ordonnées par objet, multiplicité et codomaine. Le domaine mathématique contient \(81\) positions, regroupées en
codomaines distincts. Le domaine mathématique contient \(81\) positions, regroupées en
\(56\) familles de trajectoires et \(13\) classes d'extension ; les \(324\)
trajectoires orientées ajoutent l'orientation initiale et ne représentent pas de masses
ou de particules supplémentaires. La colonne ``opérateur d'évaluation'' n'impose pas de masse là où l'objet exige
une carte nulle, une composition préalable, un invariant topologique ou une
obstruction au prolongement.

\subsection{Définitions spectrales de particule, spin, couleur et saveur}

Une \emph{particule} est une section survivante réalisée comme parcours
orienté, munie d'un enregistrement, d'une représentation de symétrie, d'un opérateur spectral
et d'une carte physique. Un fermion réalise la graduation extérieure et les relations
CAR ; un boson réalise la complétion symétrique et les relations CCR. Le
spin enregistre l'holonomie de l'orientation : le retour de \(2\pi\) clôt
la phase bosonique, tandis que la section spinorielle retrouve son état après
\(4\pi\). Cette distinction relie le double feuillet surfacique--volumétrique à la
parité et à l'exclusion de Pauli.

La \emph{couleur} est la représentation ternaire qui relève la coordonnée
\(r-c\pmod3\) vers l'orbite fondamentale puis vers l'adjointe de
\(\mathfrak{su}_3\). La \emph{saveur} identifie une base de réalisation
au sein d'une représentation de charge, couleur et spin fixés ; CKM et PMNS
assurent le transport entre la base d'interaction et la base spectrale. Les neutrinos
sont des sections fermioniques neutres : leurs étiquettes de saveur relèvent du
lecteur de mélange et leurs masses des vecteurs propres de l'opérateur correspondant.

Une particule virtuelle est une section finie avec enregistrement et propagateur jusqu'à
un horizon de prolongement ; une réalisation de Majorana est le secteur
fermionique réel fixé par conjugaison ; et les secteurs de Fibonacci, d'Ising et de
\(\mathbb Z/6\mathbb Z\) donnent des dimensions quantiques, des règles de fusion et
des caractères de tresse. Chaque définition conserve ainsi le codomaine que la
structure lui assigne.

